\chapter{Monitoring Jobs and Inflation (Text Ch. 21)}

\section{Employment and unemployment}

\subsection{Why unemployment is a problem}

Unemployment results in
\begin{itemize}
  \item \textbf{lost incomes and production}: unemployed workers lose income (only
        partly replaced by Employment Insurance), and society loses the output they would
        have produced; and
  \item \textbf{lost human capital}: a long spell of unemployment erodes skills and
        work habits and can permanently damage job prospects, health and self-esteem.
\end{itemize}

\subsection{The Labour Force Survey}

Statistics Canada measures employment and unemployment through the monthly \emph{Labour
Force Survey} (LFS), which asks about 54{,}000 households about the labour market
activity of their members.

The population is divided as follows.
\begin{itemize}
  \item The \textbf{working-age population}: the total number of people aged 15 years and
        over (excluding residents of the territories and First Nations reserves, full-time
        members of the armed forces and people in institutions). It is split into
        \begin{itemize}
          \item people \textbf{in the labour force}, and
          \item people \textbf{not in the labour force} (retirees, full-time students not
                looking for work, people caring for family, discouraged workers).
        \end{itemize}
  \item People too young to work (under 15).
\end{itemize}

\begin{definition}[Labour force, employed, unemployed]
The \emph{labour force} is the sum of the employed and the unemployed. A person is
\emph{employed} if they did any paid work (full-time or part-time) or had a job but were
absent from it. A person is \emph{unemployed} if they had no work but were available for
work and
\begin{enumerate}
  \item were on \emph{temporary layoff} with an expectation of recall, or
  \item had made specific efforts to find a job within the previous four weeks, or
  \item were waiting to start a new job within four weeks.
\end{enumerate}
\end{definition}

\begin{important}[Pitfall]
The three criteria above define who is \emph{unemployed}, not employed. Someone who
wants a job but has not looked in the last four weeks is not in the labour force at all.
\end{important}

\section{Four labour market indicators}

\begin{definition}[Unemployment rate]
The percentage of the labour force that is unemployed:
\[
  \text{unemployment rate} = \frac{\text{number of people unemployed}}{\text{labour force}}
  \times 100.
\]
\end{definition}

\begin{definition}[Involuntary part-time rate]
The percentage of the labour force who work part time but want full-time jobs:
\[
  \text{IPTR} = \frac{\text{number of involuntary part-time workers}}{\text{labour force}}
  \times 100.
\]
\end{definition}

\begin{definition}[Labour force participation rate]
The percentage of the working-age population who are members of the labour force:
\[
  \text{LFPR} = \frac{\text{labour force}}{\text{working-age population}} \times 100.
\]
\end{definition}

\begin{definition}[Employment-to-population ratio]
The percentage of the working-age population who have jobs:
\[
  \text{EPR} = \frac{\text{employment}}{\text{working-age population}} \times 100.
\]
\end{definition}

\begin{example}[Computing the indicators]
In a month (thousands of people): working-age population 30{,}000, employed 18{,}600,
unemployed 1{,}400; 900 of the employed are involuntary
part-time workers.
\begin{align*}
  \text{labour force} &= 18{,}600 + 1{,}400 = 20{,}000,\\
  \text{unemployment rate} &= 1{,}400/20{,}000 \times 100 = 7.0\%,\\
  \text{IPTR} &= 900/20{,}000 \times 100 = 4.5\%,\\
  \text{LFPR} &= 20{,}000/30{,}000 \times 100 \approx 66.7\%,\\
  \text{EPR} &= 18{,}600/30{,}000 \times 100 = 62.0\%.
\end{align*}
Now suppose 200 of the unemployed give up looking (become \emph{discouraged searchers}).
The labour force falls to $19{,}800$, unemployment to $1{,}200$, and the unemployment
rate to $1{,}200/19{,}800 \approx 6.1\%$. The unemployment rate \emph{fell}, though no one
found a job; the LFPR fell to 66.0\% and the EPR did not change.
\end{example}

\begin{remark}
A useful identity: $\text{EPR} = \text{LFPR} \times (1 - u)$, where $u$ is the
unemployment rate as a decimal. Check: $0.667 \times 0.93 \approx 0.62$.
\end{remark}

\subsection{Other definitions of unemployment}

The purpose of the unemployment rate is to measure the underutilization of labour
resources. The official measure is imperfect because it excludes:
\begin{enumerate}
  \item \emph{Discouraged searchers}: people who are neither working nor looking for work
        but are available for a job and have stopped searching because of repeated
        failure.
  \item \emph{Long-term future starts}: people with a job to start in \emph{more} than four
        weeks are not counted as unemployed.
  \item \emph{Involuntary part-timers}: part-time workers who want full-time jobs but
        can't find them are counted as employed.
\end{enumerate}
Statistics Canada publishes supplementary measures (R1 to R8) that include some of these
groups; the broadest (R8) is several points higher than the official rate.

\section{Unemployment and full employment}

There is always some unemployment. Economists classify it into three types.

\begin{definition}[Frictional unemployment]
Unemployment that arises from normal labour market turnover: people entering and
leaving the labour force and the ongoing creation and destruction of jobs. It is a
permanent and healthy phenomenon in a dynamic, growing economy (job search takes time).
\end{definition}

\begin{definition}[Structural unemployment]
Unemployment that arises when changes in technology or international competition change
the skills needed to perform jobs or change the locations of jobs. It usually lasts
longer than frictional unemployment because workers must retrain or relocate.
\end{definition}

\begin{definition}[Cyclical unemployment]
The higher-than-normal unemployment at a business cycle trough and lower-than-normal
unemployment at a peak. It fluctuates with the business cycle: it rises in recessions and
falls (can be negative) in expansions.
\end{definition}

\begin{definition}[Natural unemployment and full employment]
\emph{Natural unemployment} is the unemployment that arises from frictional and
structural unemployment when there is no cyclical unemployment.
\[
  \text{natural unemployment rate} = \frac{\text{natural unemployment}}{\text{labour force}}
  \times 100.
\]
\emph{Full employment} is the situation in which the unemployment rate equals the
natural unemployment rate. At full employment there is no cyclical unemployment.
\end{definition}

The natural unemployment rate itself changes over time. Key factors are:
\begin{enumerate}
  \item the \emph{age distribution} of the population (young workers change jobs more
        often, so a younger labour force has more frictional unemployment);
  \item the \emph{scale of structural change} (rapid technological change raises it);
  \item the \emph{real wage rate} (a minimum wage or efficiency wage above the
        market-clearing level creates a surplus of labour);
  \item \emph{unemployment benefits} (more generous EI lowers the opportunity cost of
        search and lengthens it).
\end{enumerate}

\section{Real GDP and unemployment over the cycle}

\begin{definition}[Output gap]
\emph{Potential GDP} is the quantity of real GDP produced at full employment; it is the
economy's capacity to produce on a sustained basis.
\[
  \text{output gap} = \text{real GDP} - \text{potential GDP},
\]
often expressed as a percentage of potential GDP.
\end{definition}

When the output gap is positive (an \emph{inflationary gap}), the unemployment rate is
below the natural rate; when it is negative (a \emph{recessionary gap}), unemployment
exceeds the natural rate. This inverse relationship is \emph{Okun's law}: roughly, each
percentage point of cyclical unemployment corresponds to an output gap of about 2\% of
potential GDP.

\begin{example}
Potential GDP is \$2{,}000 billion and real GDP is \$1{,}960 billion. The output gap is
$-40$ billion, or $-2\%$ of potential. By Okun's rule of thumb the unemployment rate is
about 1 percentage point above the natural rate; if the natural rate is 6\%, unemployment
is about 7\%.
\end{example}

\section{The price level, inflation and deflation}

\begin{definition}[Price level, inflation, deflation]
The \emph{price level} is the average level of prices (and so determines the value of
money). A persistently rising price level is \emph{inflation}; a persistently falling
price level is \emph{deflation}. The \emph{inflation rate} is the annual percentage
change in the price level.
\end{definition}

\subsection{Why inflation and deflation are problems}

A low, steady and anticipated inflation rate is not a big problem. An
\emph{unpredictable} (unanticipated) inflation or deflation is, because it causes
\begin{enumerate}
  \item \textbf{redistribution of income}: workers on fixed money wage contracts lose
        when inflation is higher than expected (employers gain), and vice versa;
  \item \textbf{redistribution of wealth}: unexpected inflation lowers the real value of
        debts, so borrowers gain and lenders lose (unexpected deflation does the
        opposite);
  \item \textbf{lower real GDP and employment}: uncertainty makes firms cautious, so they
        invest and hire less;
  \item \textbf{diversion of resources from production}: people spend effort forecasting
        and protecting themselves from inflation instead of producing.
\end{enumerate}

\begin{definition}[Hyperinflation]
At its worst, inflation becomes \emph{hyperinflation}: an inflation rate of 50\% a month
or higher. Money loses value so quickly that workers are paid twice a day and spend their
pay immediately (Germany 1923, Zimbabwe 2007--2008, Venezuela after 2016).
\end{definition}

\section{The Consumer Price Index}

\begin{definition}[CPI]
The \emph{Consumer Price Index} (CPI) measures the average of the prices paid by
consumers for a fixed basket of consumer goods and services. Statistics Canada calculates
it every month.
\end{definition}

\textbf{Reading the CPI numbers.} The CPI is defined to equal 100 for the
\emph{reference base period} (in Canada, 2002 $= 100$). A CPI of 137 means the basket costs
37\% more than in the base period.

\subsection{Constructing the CPI}

\begin{enumerate}
  \item \textbf{Selecting the CPI basket}: based on Statistics Canada's Survey of
        Household Spending; the weights reflect shares of spending (shelter,
        transportation and food are the largest).
  \item \textbf{Conducting the monthly price survey}: every month, Statistics Canada
        checks the prices of the basket items in cities across Canada.
  \item \textbf{Calculating the CPI}: find the cost of the basket at base-period prices,
        find the cost of the same basket at current-period prices, then
\[
  \text{CPI} = \frac{\text{cost of CPI basket at current-period prices}}
  {\text{cost of CPI basket at base-period prices}} \times 100.
\]
\end{enumerate}

\subsection{Measuring the inflation rate}

A major purpose of the CPI is to measure the inflation rate:
\[
  \text{inflation rate} = \frac{\text{CPI this year} - \text{CPI last year}}
  {\text{CPI last year}} \times 100.
\]

\begin{example}[CPI and inflation]
The basket is 10 oranges and 5 haircuts.
\begin{center}
\begin{tabular}{lccc}
\toprule
 & Base period & Year 1 & Year 2 \\
\midrule
Price of an orange & \$1.00 & \$1.20 & \$1.30 \\
Price of a haircut & \$8.00 & \$9.00 & \$9.50 \\
Cost of basket & $10 + 40 = \$50$ & $12 + 45 = \$57$ & $13 + 47.50 = \$60.50$ \\
CPI & 100 & 114 & 121 \\
\bottomrule
\end{tabular}
\end{center}
Inflation from year 1 to year 2 is $(121 - 114)/114 \times 100 \approx 6.1\%$.
\end{example}

\begin{important}[Pitfall: percentage points vs percent]
The CPI rose by 7 \emph{points}, but inflation is 6.1 \emph{percent}, because the
change must be divided by last year's level (114), not by 100. Only when last year's CPI
is 100 are the two the same.
\end{important}

\subsection{Real variables and nominal variables}

Using a price index we can convert a nominal (dollar) value into a real value, measured
in base-period dollars:
\[
  \text{real value} = \frac{\text{nominal value}}{\text{price index}} \times 100.
\]

\begin{example}[Real wage]
The average hourly wage was \$18 in 2005 (CPI 107) and \$28 in 2020 (CPI 137). In 2002
dollars the real wage was $18/1.07 \approx \$16.82$ in 2005 and $28/1.37 \approx \$20.44$
in 2020: the real wage rose by about 21.5\%, much less than the 55.6\% nominal increase.
\end{example}

The \emph{real interest rate} is the analogous idea for interest rates (see the finance
chapter): real interest rate $\approx$ nominal interest rate $-$ inflation rate.

\subsection{The biased CPI}

The CPI may \emph{overstate} the true inflation rate for four reasons:
\begin{enumerate}
  \item \textbf{New goods bias}: new goods (e.g.\ smartphones) replace older goods; if the
        new ones are more expensive but better, the CPI records a price rise that is really
        an improvement.
  \item \textbf{Quality change bias}: products improve every year (safer cars, faster
        computers); part of any price rise pays for higher quality, not inflation.
  \item \textbf{Commodity substitution bias}: when a good's relative price rises, consumers
        switch to substitutes, but the fixed basket assumes they keep buying the same
        quantities.
  \item \textbf{Outlet substitution bias}: consumers shift to cheaper retailers (discount
        stores, online shopping), which the CPI does not fully capture.
\end{enumerate}
Estimates suggest the Canadian CPI overstates inflation by roughly 0.5 to 0.6 percentage
points a year. The bias distorts private contracts and raises indexed government outlays
(CPP and Old Age Security are indexed to the CPI).

\subsection{Alternative price indexes}

\begin{itemize}
  \item \textbf{GDP deflator}: a broader measure, covering consumption, investment,
        government expenditure and net exports:
        \[
          \text{GDP deflator} = \frac{\text{nominal GDP}}{\text{real GDP}} \times 100.
        \]
  \item \textbf{Chained price index for consumption} (CPIC): an index of the prices of all
        items in consumption expenditure in GDP, using current weights (so it has less
        substitution bias):
        \[
          \text{CPIC} = \frac{\text{nominal consumption expenditure}}
          {\text{real consumption expenditure}} \times 100.
        \]
\end{itemize}

\subsection{Core inflation}

\begin{definition}[Core inflation]
The \emph{core inflation rate} excludes the most volatile elements of the CPI (such as
food and energy) to reveal the underlying trend. The Bank of Canada historically used
CPIX (excluding the eight most volatile components); since 2016 it focuses on CPI-trim,
CPI-median and CPI-common.
\end{definition}

\section{Summary}

\begin{itemize}
  \item The LFS classifies the working-age population (15+) as employed, unemployed, or
        not in the labour force. Labour force $=$ employed $+$ unemployed.
  \item Indicators: unemployment rate, IPTR (both over the labour force), LFPR and EPR
        (both over the working-age population). Discouraged searchers lower the measured
        unemployment rate.
  \item Unemployment is frictional, structural or cyclical; natural unemployment $=$
        frictional $+$ structural; full employment means zero cyclical unemployment.
  \item Output gap $=$ real GDP $-$ potential GDP; unemployment above the natural rate
        goes with a negative gap (Okun's law).
  \item CPI $=$ cost of basket at current prices $/$ cost at base prices $\times 100$;
        inflation is the percentage change in the CPI. The CPI overstates inflation (new
        goods, quality, commodity and outlet substitution).
  \item Real value $=$ nominal value $/$ price index $\times 100$. Alternatives: GDP
        deflator, CPIC, core inflation.
\end{itemize}
